\section{Контентная фильтрация}

В данном разделе рассмотрены методы рекомендаций рецептов, учитывающие их содержание.
С помощью оценок пользователя, для которого необходимо сделать рекомендацию, и вектора параметров рецепта предсказывается, как данный пользователь будет оценивать рецепты.

\subsection{Параметры рецепта}
Чтобы работать с содержанием рецептов, необходимо формализовать его.
Каждый рецепт должен быть представлен как вектор вещественных чисел.
Причём соответствующие координаты векторов должны описывать один и тот же параметр.
Таким образом, необходимо описать способ получения числовых параметров рецепта.

\subsubsection{Преобразование текста в вектор}
Рецепты представляются текстовой информацией, в которую входит:
\begin{itemize}
    \item название рецепта;
    \item описание рецепта;
    \item описание шагов;
    \item список ингредиентов;
    \item список тегов при их наличии.
\end{itemize}
Вводится определение документа и словаря.
Документ -- список всех слов из всех перечисленных компонентов рецепта, а словарь -- список рассматриваемых слов~\cite{segaran}.
Пусть:
\begin{itemize}
    \item $N$ -- количество рецептов;
    \item $D_i$, $i \in \overline{1,N}$ -- документ номер $i$;
    \item $M$ -- количество рассматриваемых слов;
    \item $W_j$, $j \in \overline{1,M}$ -- слово номер $j$ в словаре;
    \item $\operatorname{count}(W_j, D_i)$ -- количество вхождений слова $W_j$ в документ $D_i$;
    \item $\operatorname{docs}(W_j)$ -- количество документов, содержащих слово $W_j$;
    \item $X$ -- матрица размера $N \times M$, где строка $i$ является представлением рецепта $i$.
\end{itemize}
Значения $X_{ij}$ рассчитывается по формуле~\eqref{eq:tfidf}:
\begin{equation}
    X_{ij} = \frac{\text{count}(W_j, D_i)}{|D_i|} \cdot \ln \left( \frac{N}{\text{docs}(W_j)} \right).
    \label{eq:tfidf}
\end{equation}
Таким образом определяется набор векторов, представляющих рецепты~\cite{salton1975vector}.
Данные вектора могут быть дополнены числовыми параметрами рецепта, такими как:
\begin{itemize}
    \item количество времени, требуемое на выполнение рецепта;
    \item количество ингредиентов в рецепте;
    \item количество шагов рецепта;
    \item параметры <<пищевой ценности>>.
\end{itemize}
Если какой-либо из перечисленных параметров используется, то необходимо применить нормализацию~\cite{sharma2022scaling}.

\subsection{Классификация Байеса}
Существует модель предсказания рейтинга объекта, основанная на формуле Байеса~\cite{aggarwal}.
Вводятся следующие обозначения:
\begin{itemize}
    \item $U$ -- множество из $n$ векторов размерности $m$, содержащих параметры объектов, оценённых пользователем;
    \item $T_i$ -- оценка пользователем объекта $i, i \in \overline{1,N}$;
    \item $X$ -- вектор параметров объекта, для которого необходимо предсказать оценку;
    \item $f(X)$ -- функция оценки пользователя для объекта $X$.
\end{itemize}
Классификатор Байеса предполагает, что оценка пользователем объекта дискретна, а также параметры объектов принимают значения $0$ и $1$~\cite{aggarwal}.
Для более точных предсказаний рекомендуется, чтобы оценка принимала лишь два возможных значения, но в общем случае оценка $\alpha \in A$.
Классификация основана на оценке вероятности $P\{f(X) = \alpha | X = (x_1,\dots,x_m)\}$.
По формуле Байеса~\eqref{eq:bias_init}:
\begin{equation}
    \label{eq:bias_init}
    \small P\{f(X) = \alpha | X = (x_1,\dots,x_m)\} = \frac{P\{f(X)=\alpha\} \cdot P\{X=(x_1,\dots,x_m)|f(X)=\alpha\}}{P\{X=(x_1,\dots,x_m)\}}.
\end{equation}
Делается допущение, что составляющие вектора $X$ принимают своё значение независимо друг от друга~\cite{aggarwal}.
Следовательно, справедливо равенство~\eqref{eq:p_decompose}:
\begin{equation}
    \label{eq:p_decompose}
    P\{X=(x_1,\dots,x_m)|f(X)=\alpha\}=\prod_{i=1}^{m}{P\{X_i = x_i|f(X)=\alpha\}}.
\end{equation}

Значение вероятности $P\{X=(x_1,\dots,x_m)\}$ определяется по формуле полной вероятности~\eqref{eq:full_prob}.
Учитывая~\eqref{eq:p_decompose}, она принимает вид~\cite{aggarwal}:
\begin{equation}
\label{eq:full_prob}
    \small P\{X=(x_1,\dots,x_n)\}=\\
    \sum_{\alpha \in A}{P\{f(X)=\alpha\} \cdot \prod_{i=1}^{m}{P\{X_i=x_i|f(X)=\alpha\}}}.
 \end{equation}
Пусть $U_{\alpha} = \{u_i | u_i \in U, T_i = \alpha\}$.
Тогда вероятность $P\{f(X)=\alpha\}$ оценивается по формуле~\eqref{eq:prob1}:
\begin{equation}
\label{eq:prob1}
    P\{f(X)=\alpha\}=\frac{|U_{\alpha}| + \delta}{|U| + 2\delta},
\end{equation}
где $\delta$ -- малая положительная величина, называемая параметром сглаживания~\cite{aggarwal}.
Пусть $U^{(x_i)}_{\alpha} = \{u | u \in U_{\alpha},u_i=x_i\}$.
Тогда для оценки значения $P\{X_i=x_i|f(X)=\alpha\}$ используется формула~\eqref{eq:prob2}:
\begin{equation}
\label{eq:prob2}
    P\{X_i=x_i|f(X)=\alpha\}=\frac{|U^{(x_i)}_{\alpha}| + \gamma}{|U_{\alpha}| + 2\gamma},
\end{equation}
где $\gamma$ -- параметр сглаживания.
В результате формула~\eqref{eq:bias_init} может быть вычислена $\forall \alpha \in A$.
Тогда оценка объекта вычисляется как~\eqref{eq:res_score_Bias}:
\begin{equation}
\label{eq:res_score_Bias}
    \hat{\alpha} = \sum_{\alpha \in A}{P\{f(x)=\alpha|x=X\}} \cdot \alpha.
\end{equation}

\subsection{Регрессия}
Предполагается, что зависимость целевого параметра от набора числовых признаков имеет определённый вид.
Таким образом, задаётся формула, содержащая несколько параметров, которая устанавливает связь между параметрами объекта и его оценкой.
Задача обучения регрессии заключается в поиске этих параметров.
Для определения качества регрессии используется метрика <<MSE>>~\cite{aggarwal}, вычисляемая по формуле~\eqref{eq:mse}:
\begin{equation}
	\label{eq:mse}
	\operatorname{Loss} = \frac{1}{n}\sum_{i=1}^{n}{(y_i - \hat{y}_i)^2},
\end{equation}
где:
\begin{itemize}
	\item $y_i$ -- эталонное значение результата модели;
	\item $\hat{y}_i$ -- результат модели;
	\item $n$ -- количество тестовых случаев.
\end{itemize}
Чем меньше значение данной метрики, тем предпочтительнее модель по сравнению с другими.

\subsubsection{Линейная регрессия}
Рассмотрена линейная регрессия.
Для описания данного метода вводятся следующие обозначения:
\begin{itemize}
    \item $X$ -- матрица $n \times m$, строка $i$ которой содержит параметры объекта с номером $i$;
    \item $w$ -- вектор из $m$ весов параметров объектов;
    \item $y$ -- вектор оценок объектов пользователем.
\end{itemize}
Предполагается, что зависимость оценки от признаков рецепта линейна и имеет вид~\eqref{eq:linear_reg}:
\begin{equation}
	\label{eq:linear_reg}
	y = a_0 + a_1x_1 + a_2x_2 + . . . + a_mx_m,
\end{equation}
где:
\begin{itemize}
	\item $m$ -- количество признаков объекта;
	\item $a_i$ -- параметры модели;
	\item $x_i$ -- признаки объектов.
\end{itemize}

Вектор весов $w$ вычисляется по формуле~\eqref{eq:weights}:
\begin{equation}
\label{eq:weights}
    w^T = (X^TX)^{-1}X^Ty.
\end{equation}
Линейная регрессия подходит для случаев, когда рейтинг описывается числом из определённого диапазона.

\subsection{Регуляризация линейной регрессии}
Задача регрессии не является устойчивой~\cite{aggarwal}.
Это значит, что найденные параметры модели могут оказаться такими, что она хорошо работает на одном наборе данных, но плохо на другом.
Данная проблема решается с помощью регуляризации.
Функция потерь <<MSE>>, которая оптимизируется \\ \\в ходе обучения модели, модифицируется следующим образом~\eqref{eq:mse_l2}:
\begin{equation}
	\label{eq:mse_l2}
	\operatorname{Loss} = \frac{1}{n}\sum_{i=1}^{n}{(y_i - \hat{y}_i)^2} + \alpha \cdot \sum_{j=1}^{p}{w_j^2},
\end{equation}
где:
\begin{itemize}
	\item $p$ -- количество параметров (весов);
	\item $w$ -- параметры (веса);
	\item $\alpha$ -- коэффициент регуляризации.
\end{itemize}
В результате регуляризации формула расчёта весов $w$~\eqref{eq:weights_reg} принимает вид:
\begin{equation}
\label{eq:weights_reg}
	 w^T = (X^TX + \alpha I)^{-1}X^Ty,
\end{equation}
где $I$ -- единичная матрица размером $n \times n$.

\subsubsection{Нормализация регрессии}
Для минимизации функции потерь в некоторых случаях следует провести нормализацию матрицы $X$~\cite{aggarwal}.
Она производится в соответствии с формулами~\eqref{eq:min_max_norm}--\eqref{eq:std_norm}.
Но наиболее подходящий метод для линейной регуляризованной регрессии основан на формуле~\eqref{eq:std_norm}~\cite{aggarwal}.
Далее представлены формулы нормализации:
\begin{equation}
	\label{eq:min_max_norm}
	x'_{ij}=\frac{x_{ij}-\operatorname{min}(X_j)}{\operatorname{max}(X_j) - \operatorname{min}(X_j)},
\end{equation}
где:
\begin{itemize}
	\item $x_{ij}$ и $x'_{ij}$ -- начальное и нормализованное значение признака $j$ объекта $i$ соответственно;
	\item $\operatorname{min}(X_j)$ -- минимальное значение признака $j$;
	\item $\operatorname{max}(X_j)$ -- максимальное значение признака $j$.
\end{itemize}
\begin{equation}
	\label{eq:unit_norm}
	x_i = \frac{x_i}{\lVert x_i \rVert},
\end{equation}
где $x_i$ -- вектор объекта $i$.
\begin{equation}
	\label{eq:std_norm}
	x'_{ij}=\frac{x_{ij}-\mu_j}{\sigma_j},
\end{equation}
где:
\begin{itemize}
	\item $x_{ij}$ и $x'_{ij}$ -- начальное и нормализованное значение признака $j$ объекта $i$ соответственно;
	\item $\mu_j$ -- среднее арифметическое значение для признака $j$;
	\item $\sigma_j$ -- стандартное отклонение для признака $j$.
\end{itemize}
В рекомендательной системе нормализация всех векторов объектов проводится заранее, чтобы избежать дополнительных временных затрат, так как набор векторов фиксированный.

\subsubsection{Другие модели регрессии}
Для разных типов оценок пользователя следует применять разные модели регрессии.
В таблице~\ref{tab:models} приведены рекомендации по выбору регрессии для различных форматов оценок объектов~\cite{aggarwal}.

\begin{table}[h]
\centering
\caption{\centering Сопоставление модели регрессии и типа оценки объектов}
\begin{tabular}{|l|l|}
\hline
\textbf{Модель регрессии} & \textbf{Тип оценки} \\
\hline
Линейная, полиномиальная и ядерная регрессии & Вещественная оценка \\
\hline
Многомерная логистическая регрессия & Дискретная оценка \\
\hline
Логистическая или пробит регрессия & Бинарная оценка \\
\hline
\end{tabular}
\label{tab:models}
\end{table}
Таким образом, регрессию можно использовать для любых типов оценок.

\subsubsection{Метод уменьшения размерности}
Главная проблема регрессии заключается в том, что размер вектора, представляющего рекомендуемый объект, больше, чем количество оценок пользователя.
В этом случае регрессия не может обучиться на предоставленных ей данных так, чтобы выдавать результат точнее случайного угадывания~\cite{pazzani2007content}.
Пусть:
\begin{itemize}
	\item $N$ -- количество возможных оценок объекта пользователем;
	\item $a$ -- объект, который необходимо оценить;
	\item $f(a, b)$ -- метрика расстояния между $a$ и $b$.
\end{itemize}
Для решения описанной проблемы выполняются следующие шаги~\cite{pazzani2007content}:
\begin{itemize}
	\item получается множество векторов объектов $X_s$, оценённых пользователем на $s$, $\forall s \in \overline{1, N}$;
	\item вычисляется средний вектор $\omega_s$ для каждого $X_s$, $s \in \overline{1, N}$;
	\item если вектор $|X_s| = 0$, то вектор $\omega_s$ принимается нулевым.
	\item вычисляется $d_s = f(a, \omega_s)$ для каждого $s  \in \overline{1, N}$;
	\item вектор $y = (d_1, d_2, \dots, d_N)$ используется для обучения регрессионной модели.
\end{itemize}
Таким образом, алгоритм на основе регрессии даёт более точные рекомендации, так как он способен обучиться на меньшем количестве оценок.
Множество $\{X_s | s \in \overline{1, N}\}$ называется профилем пользователя.
Чтобы применить полученный метод предсказания оценки к объекту, сначала вычисляется вектор $y$ с помощью профиля по описанному выше алгоритму, а затем к $y$ применяется регрессия.
В результате на выходе регрессии получается предсказанная оценка пользователя для рассматриваемого объекта.
В качестве метрики $f$ используется косинусная близость~\eqref{eq:cosinesim} или Евклидова расстояние~\eqref{eq:dist}.
\begin{equation}
	f(a, b) = \frac{a \cdot b}{\lVert a\rVert \cdot \lVert b\rVert}.
	\label{eq:cosinesim}
\end{equation}

\subsection{Метод ближайших соседей}
Вводится метрика схожести объектов $\rho(a, b)$, которая тем больше, чем более разные $a$ и $b$.
Необходимо предсказать оценку для $x$, основываясь на оценках $y_i, i \in \overline{1,n}$ для множества объектов $U = \{u_i, i \in \overline{1,n}\}$, которые оценил пользователь.
С помощью метрики $\rho$ находится ровно $k$  элементов множества $U$, которые наиболее близки к $x$~\cite{aggarwal}.
Пусть они обозначаются как $u_1,\dots,u_k$.
Тогда оценка для $x$ вычисляется по формуле~\eqref{eq:knn_result}:
\begin{equation}
\label{eq:knn_result}
    \hat{\alpha} = \frac{\sum\limits_{i=1}^{k} \rho(x,u_i) y_i}{\sum\limits_{i=1}^{k} \rho(x,u_i)}.
\end{equation}
Особенностью данного метода является необходимость вычислять метрику для $x$ и каждого элемента множества $U$.
Для оптимизации алгоритм <<$k$ ближайших соседей>>, используется кластеризация~\cite{aggarwal}.
Она позволяет разделить множество $U$ на $p$ групп $g_1,\dots,g_p$.
Пусть количество возможные оценки равны $1,\dots,s$.
Каждая группа $g_i, i \in \overline{1,p}$ делится на $g_{(i,1)},\dots,g_{(i,s)}$ согласно оценкам~\cite{aggarwal}.
В группу $g_{(i,j)}$ входят объекты из $g_i$, оценённые пользователем на $j$.
Тогда максимальное общее количество групп равно $p \times s$.
Каждая группа $g_{(i,j)}$ заменяется на усреднённый по группе элемент $e_{(i,j)}$, параметры которого, равны среднему из соответствующих параметров в пределах данной группы.
Оценка такого усреднённого элемента $e_{(i,j)}$ считается равной $j$~\cite{aggarwal}.
В результате алгоритм <<$k$ ближайших соседей>> применяется не к множеству $U$, а к множеству объектов $\{e_{(i,j)}\}, i \in \overline{1,p}, j \in \overline{1,s}$.
Таким образом, множество $\{e_{(i,j)}\}$ вычисляется лишь $1$ раз, но используется всякий раз для применения метода <<$k$ ближайших соседей>> к какому-либо объекту $x$~\cite{aggarwal}.
Пересчёт $\{e_{(i,j)}\}$ выполняется только при значительном пополнении выборки оценённых объектов $U$.
В качестве метрики $\rho$ достаточно взять Евклидово расстояние или формулу~\eqref{eq:rho_a_b}:
\begin{equation}
\label{eq:rho_a_b}
    \rho(a,b) = 1 - \frac{\sum\limits_{i=1}^{m} x_i y_i}{\sqrt{\sum\limits_{i=1}^{m} x_i^2 \cdot \sum\limits_{i=1}^{m} y_i^2}}.
\end{equation}

\subsection{Ядерная регрессия}
Рассматривается набор объектов, представленных в виде векторов чисел.
Каждый вектор задаёт координаты точки в многомерном пространстве.
Ядерная регрессия вычисляет неизвестное значение оценки данной точки как взвешенное среднее оценок окружающих её точек.
Алгоритм предсказания выполняется следующим образом:
\begin{itemize}
	\item для новой точки определяются расстояния до всех известных точек  из обучающей выборки;
	\item расстояния передаются в математическую функцию $K$ -- ядро, которая преобразует их в вес;
	\item итоговая оценка вычисляется как среднее значение всех оценок в выборке, умноженных на их веса.
\end{itemize}
Таким образом, применяется следующая формула~\eqref{eq:kernel_reg}:
\begin{equation}
\label{eq:kernel_reg}
	\hat{y}(x) = \frac{\sum\limits_{i=1}^{n} K \Big( \frac{x-x_{i}}{h} \Big) y_{i}}{\sum\limits_{i=1}^{n} K \Big( \frac{x-x_{i}}{h} \Big)}.
\end{equation}

\subsection{Кластеризация}
Задача кластеризации заключается в разбиении множества объектов на группы в соответствии с его структурой~\cite{segaran}.
Это значит, что в пределах одной группы объекты должны отличаться друг от друга как можно меньше.
В данной работе рассмотрен метод кластеризации методом $k$ средних.
Пусть $X$ -- матрица $n \times m$, строка $i$ которой содержит параметры объекта номер $i$, а $\rho$ -- метрика схожести двух векторов.
В качестве метрики достаточно взять Евклидово расстояние~\cite{segaran}.
Алгоритм кластеризации объектов содержит следующие шаги~\cite{segaran}:
\begin{itemize}
    \item создать $k$ случайных векторов $a_1,\dots,a_k$ размерности $m$;
    \item создать пустые множества $A_1,\dots,A_k$;
    \item для каждой строки $x_i$ матрицы $X$ с помощью $\rho$ определить ближайший вектор $a_j, j \in \overline{1,k}$ и добавить его к множеству $A_j$;
    \item для $j \in \overline{1,k}$ вычислить средний по множеству $A_j$ вектор $\alpha_j$;
    \item для $j \in \overline{1,k}$ присвоить $a_j$ значение $\alpha_j$;
    \item если набор векторов $a_1,\dots,a_k$ поменялся, то перейти к шагу (2);
    \item множества $A_1,\dots,A_k$ являются найденным разбиением.
\end{itemize}
Таким образом, исходное множество объектов распадается на ровно $k$ подмножеств, образующих группы.

\subsection{Комбинации моделей}
Для того, чтобы повысить точность предсказаний применяется комбинирование алгоритмов машинного обучения в сложную модель~\cite{dietterich2000ensemble}.
В данном методе есть базовые модели, обучаемые на исходных данных и мета-модель, объединяющая их результаты в финальный ответ.
Базовая модель работает с векторами объектов и предсказывает их оценки.
Мета-модель принимает на вход вектор результатов предсказаний базовых моделей и вычисляет итоговую оценку объекта.
Таким образом, в некоторых случаях за счёт взаимной компенсации ошибок базовых моделей результат получается более точным~\cite{dietterich2000ensemble}.

Базовые модели обучаются на одних и тех же входных данных~\cite{dietterich2000ensemble}.
Для обучения мета модели необходимо подать на вход базовым моделям данные, на которых они не обучались, и полученные предсказания передать мета-модели для её обучения.
Для того, чтобы повысить точность сложной модели, базовые модели должны выдавать ошибки на разных объектах~\cite{dietterich2000ensemble}.